\clearpage
\phantomsection% 
\addcontentsline{toc}{chapter}{KATA PENGANTAR}
%\thispagestyle{fancy}

\begin{justifying}
	\large \bfseries \centering \MakeUppercase{Kata Pengantar}\linebreak
	
	\normalsize \normalfont \justifying
	% \textit{Pada halaman ini mahasiswa berkesempatan untuk menyatakan terima kasih secara tertulis kepada pembimbing dan pihak lain yang telah memberi bimbingan, nasihat, saran dan kritik, kepada mereka yang telah membantu melakukan penelitian, kepada perorangan atau lembaga yang telah memberi bantuan keuangan, materi dan/atau sarana. Cara menulis kata pengantar beraneka ragam, tetapi hendaknya menggunakan kalimat yang baku. Ucapan terima kasih agar dibuat tidak berlebihan dan dibatasi pada pihak yang terkait secara ilmiah (berhubungan dengan subjek/materi penelitian). } \par
	Puji syukur kehadirat Allah SWT atas segala limpahan rahmat, karunia, serta hidayah-Nya sehingga penulis dapat menyelesaikan penyusunan tugas akhir ini dengan baik. Shalawat serta salam semoga senantiasa tercurah kepada Nabi Muhammad SAW yang telah membawa umat manusia menuju jalan yang penuh ilmu pengetahuan. Penyusunan tugas akhir ini tidak terlepas dari berbagai arahan, bantuan, doa, serta dukungan dari banyak pihak. Oleh karena itu, pada kesempatan ini penulis menyampaikan rasa terima kasih yang sebesar-besarnya kepada: \par
	\begin{enumerate}
		\item Kedua orang tua saya, Bapak Pamuji Waskito Raharjo dan Mamak Anik Purwanti, yang selalu memberikan dukungan, doa yang tiada henti, biaya pendidikan, serta cinta yang tulus sehingga saya dapat menyelesaikan setiap langkah dalam perjalanan perkuliahan ini.
		\item Kedua adik saya, Delfy Celcilia dan Atha Rasyid Risky, yang selalu memberikan semangat, dukungan, dan doa dalam setiap langkah perjalanan saya.
		\item Keluarga besar saya, Nenek Sarmi, Mbah Surati, keluarga besar Bapak, keluarga besar Ibu, serta seluruh kerabat yang selalu memberikan dukungan dan doa yang tulus.
		\item Bapak Martin Clinton Tosima Manullang, Ph.D. dan Ibu Leslie Anggraini, S.Kom., M.Cs. selaku Dosen Pembimbing yang telah meluangkan waktu, memberikan bimbingan, arahan, dukungan moral dan material, serta berbagai masukan yang sangat berharga selama proses penyusunan tugas akhir ini.
		\item Bapak Imam Eko Wicaksono, S.Si., M.Si., Bapak Meida Cahyo Untoro, S.Kom., M.Kom. dan Bapak Andika Setiawan, S.Kom., M.Cs. selaku Dosen Penguji yang telah memberikan saran, kritik, dan masukan yang konstruktif untuk meningkatkan kualitas tugas akhir ini.
		\item Ibu Winda Yulita, M.Cs. dan Ibu Leslie Anggraini, S.Kom., M.Cs. selaku Dosen Wali yang memberikan bimbingan, dukungan, serta perhatian selama masa perkuliahan.
		\item Bapak Hadi, Bapak Ahmad, Bapak Paul, Bapak Mesa, dan Bapak Wiko dari PT Timah Tbk yang telah memberikan ide, saran, waktu, serta menyediakan segala kebutuhan penelitian sehingga penelitian ini dapat terlaksana dengan baik.
		\item Rahma Nur Fadillah, wanita hebat yang telah menemani perjalanan perkuliahan dari awal hingga akhir, memberikan dukungan, semangat, dan doa yang tiada henti. Terima kasih atas segala bantuan dan motivasi yang selalu kamu berikan.
		\item Rekan seperjuangan Coconut Boys, INLINEER 22, Keprofesian HMIF, Grub bimbingan iMUL dan LIA, yang telah memberikan dukungan, semangat, serta berbagai bantuan selama masa perkuliahan. Terima kasih atas kebersamaan dan pertemanan yang luar biasa.
		\item Semua pihak yang tidak dapat saya sebutkan satu per satu yang telah memberikan dukungan, bantuan, serta doa selama perjalanan perkuliahan ini. Terima kasih atas segala kontribusi yang telah diberikan.
	\end{enumerate} \par
	Penulis menyadari bahwa penulisan tugas akhir ini masih jauh dari sempurna dan memiliki banyak kekurangan. Oleh karena itu, penulis sangat terbuka dan mengharapkan saran serta kritik yang bersifat membangun untuk perbaikan di masa mendatang. Akhir kata, penulis berharap semoga tugas akhir ini dapat memberikan manfaat bagi kita semua.
	\vfill
	
\end{justifying}
\clearpage





\begin{comment}
\clearpage
\pagestyle{fancy}
\fancyhf{}
\fancyhead[R]{\thepage}
\phantomsection% 
%\clearpage
%\phantomsection% 
\addcontentsline{toc}{chapter}{KATA PENGANTAR}
%\thispagestyle{fancy}

\begin{justifying}
	\large \bfseries \centering \MakeUppercase{Kata Pengantar}\linebreak
	
	\normalsize \normalfont \justifying
	Puji syukur kehadirat Tuhan Yang Maha Esa atas limpahan rahmat, karunia, serta petunjuk-Nya sehingga penyusunan tugas akhir ini telah terselesaikan dengan baik. Dalam penyusunan tugas akhir ini penulis telah banyak mendapatkan arahan, bantuan, serta dukungan dari berbagai pihak. Oleh karena itu pada kesempatan ini penulis mengucapan terima kasih kepada: \par
	\begin{enumerate}
		\item Bapak Prof. Dr. I. Nyoman Pugeg Aryantha selaku Rektor Institut Teknologi Sumatera.  
		\item Bapak Hadi Teguh Yudistira, S.T., Ph.D. selaku Dekan Fakultas Teknologi Industri.
		\item Bapak Andika Setiawan, S. Kom., M. Cs. selaku Ketua Program Studi Teknik Informatika.
		\item Bapak Ilham Firman Ashari, S. Kom., M.T. selaku Koordinator Tugas Akhir Program Studi Teknik Informatika.  
		\item Bapak Martin C. T. Manullang, Ph.D. selaku Dosen Pembimbing atas ide, waktu, tenaga, perhatian, dan masukan yang telah disumbangsihkan kepada penulis.
            \item Bapak Andika Setiawan, S. Kom., M. Cs dan Bapak Eko Dwi Nugroho, S.Kom., M.Cs. selaku Dosen Penguji atas saran dan masukan yang diberikan. 
		\item Kedua Orang Tua dan Adik yang selalu memberikan dukungan dan doa sehingga penulis dapat menyelesaikan tugas akhir ini. Kelembutan, dukungan, dan cinta yang kalian berikan selalu menjadi sumber inspirasi dan kekuatan.
		\item Teman-teman penulis yang membantu selama masa perkuliahan yang tidak bisa disebutkan satu persatu.
	\end{enumerate} \par
	Akhir kata penulis berharap semoga tugas akhir ini dapat memberikan manfaat bagi kita semua. Penulis menyadari bahwa tugas akhir ini tidak luput dari kekurangan dan kelemahan, dan penulis  terbuka untuk menerima saran, kritik, dan masukan yang membangun.
	\vfill
	
\end{justifying}
\clearpage
\end{comment}